\chapter{Conclusione}

Il lavoro introduce un tool per l'individuazione automatica di pazienti virtuali
in reti biochimiche, basato su tre elementi principali: 
la conversione automatica di modelli qualitativi (Reactome) 
in modelli quantitativi SBML, la formulazione del 
problema di stima parametrica come ottimizzazione vincolata black-box,
 e l'integrazione di diversi algoritmi evolutivi e numerici 
 (Nevergrad, Nomad, OpenAI Evolution Strategies).

I risultati sperimentali hanno dimostrato che il sistema è in grado di:
\begin{enumerate}
    \item 	garantire la stabilità dei modelli tramite vincoli espliciti su parametri e concentrazioni,
	\item 	valutare la plausibilità di ordinamenti normalizzati delle specie,
	\item 	individuare configurazioni parametriche compatibili con i vincoli biologici.
\end{enumerate}

La validazione su pathway eterogenei ha evidenziato 
la capacità del framework di discriminare ordinamenti 
non realizzabili e di generare configurazioni coerenti 
con fenomeni biologici complessi, mantenendo tempi e uso di memoria ragionevoli.

